\section*{Hoofdstuk \ref{ch:b2:magmatic-series} --- Magmatische reeksen en tektonische omgevingen}

\begin{solution}{exo:b2:magmatic-series:1}
(a) $\ce{SiO2} = 49.60$, alkaliën $2.22 + 0.51 = 2.73$: tussen 45 en 52 en onder 5,
\emph{basalt} (B). (b) 58.5 en $3.9 + 1.6 = 5.5$: tussen 57 en 63; de bovenste
grens loopt daar van (57, 5.9) naar (63, 7.0), en ligt bij 58.5 op $5.9 + 1.1\times1.5/6 = 6.18$;
5.5 ligt eronder: \emph{andesiet} (O2). (c) 72.0 en $4.0 + 4.2 = 8.2$:
voorbij de verticaal bij 69 boven 8 en boven de lijn van (69, 8) naar (77, 0), die
bij 72 op 5.0 ligt: \emph{rhyoliet} (R). (d) 47.0 en $4.5 + 2.5 = 7.0$: bij 47 ligt het veld U1
tussen de lijn van (45, 5) naar (49.4, 7.3), op $5 + 2.3\times2/4.4 = 6.05$, en
de lijn van (45, 9.4) naar (49.4, 7.3), op $9.4 - 2.1\times2/4.4 = 8.45$. Het gesteente ligt
in U1: een \emph{tefriet} als zijn norm minder dan \qty{10}{\percent} olivijn bevat, een
\emph{basaniet} als het meer is.
\end{solution}

\begin{solution}{exo:b2:magmatic-series:2}
(a) De oxiden tellen op tot 98.40 en het totaal met vluchtige bestanddelen tot 100.00;
\ce{H2O} $1.40 < 2$ en \ce{CO2} $0.20 < 0.5$: de analyse is aanvaardbaar. (b) Factor
$100/98.40 = 1.01626$: \ce{SiO2} 52.44, \ce{TiO2} 1.22, \ce{Al2O3} 16.36,
\ce{FeO_T} 9.96, \ce{MgO} 6.71, \ce{CaO} 8.33, \ce{Na2O} 3.56, \ce{K2O} 1.42
(som 100.00); alkaliën $3.56 + 1.42 = 4.98$. (c) Ervoor: 51.60 en 4.90, net onder de
grens basalt--basaltisch andesiet bij 52: \emph{basalt}. Erna: 52.44 en 4.98; bij
52.44 ligt de bovengrens van O1, van (52, 5) naar (57, 5.9), op $5 + 0.9\times0.44/5
= 5.08$: \emph{basaltisch andesiet} (O1). De herberekening bracht het gesteente over een
grens: namen dicht bij een lijn hangen ervan af. (d) Bij 52.44 ligt de laag-K-lijn op $0.3 +
0.045\times4.44 = 0.50$ en de hoog-K-lijn op $1.2 + 0.095\times4.44 = 1.62$; 1.42
ligt ertussen: \emph{middel-K}.
\end{solution}

\begin{solution}{exo:b2:magmatic-series:3}
(a) $Q + A + P = 90$: $Q = 25/90 = \qty{27.8}{\percent}$ en $100\,P/(A + P) =
100\times35/65 = 53.8$: veld 3b, een \emph{monzograniet}, een graniet. Biotiet en
hoornblende tellen niet mee in de QAPF-som. (b) $Q + A + P = 65$: $Q = 5/65 =
\qty{7.7}{\percent}$ (tussen 5 en 20) en $100\times55/60 = 91.7$ (boven 90): veld
10*, kwartsdioriet, kwartsgabbro of kwartsanorthosiet. Het \omterm{def:b2:magmatic-series:qapf}{QAPF-diagram} alleen kiest
niet tussen hen; het plagioklaas en de mafische mineralen (hier
\qty{35}{\percent} clinopyroxeen) doen dat wel.
\end{solution}

\begin{solution}{exo:b2:magmatic-series:4}
De laag-K-lijn is $0.3 + 0.045(\ce{SiO2} - 48)$, de hoog-K-lijn $1.2 +
0.095(\ce{SiO2} - 48)$. (a) Bij 50: 0.39 en 1.39; $0.40$ ligt net boven de onderste lijn:
\emph{middel-K}. (b) Bij 56: 0.66 en 1.96: \emph{middel-K}. (c) Bij 62: 0.93 en 2.53;
$2.90 > 2.53$: \emph{hoog-K}. (d) Bij 66: 1.11 en 2.91; $1.00 < 1.11$: \emph{laag-K}.
\end{solution}

\begin{solution}{exo:b2:magmatic-series:5}
(a) Basalt: $A + F + M = 3.2 + 14.5 + 4.5 = 22.2$; $A = 14.4$, $F = 65.3$, $M = 20.3$.
Andesiet: $5.05 + 6.1 + 3.8 = 14.95$; $A = 33.8$, $F = 40.8$, $M = 25.4$. (b) De lijn
komt nooit voorbij $F = 54.5$; de basalt, met $F = 65.3$, ligt aan de kant van F
wat zijn $A$ ook is: \emph{tholeïtisch}, een ferrobasalt van een Fennertrend. Voor de
andesiet ligt de lijn bij $A = 33.8$ tussen (49.6, 35.0) en (52.5, 30.0), op $F =
49.6 + 2.9\times1.2/5 = 50.3$; $40.8 < 50.3$: \emph{kalkalkalisch}. (c) Magnetiet,
dat vanaf het begin ijzer(III) uit een geoxideerd \omterm{def:b2:magma-emplacement:magma}{magma} afvoert en verhindert dat $F/M$
even snel stijgt als in een tholeïet.
\end{solution}

\begin{solution}{exo:b2:magmatic-series:6}
MORB: $\mathrm{Nb_N} = 4/0.658 = 6.08$, $\mathrm{La_N} = 4.12/0.648 = 6.36$, verhouding
0.956. \omterm{def:b2:elements-in-the-earth:morb-oib}{Oceaaneilandbasalt}: $38/0.658 = 57.8$, $30.8/0.648 = 47.5$, verhouding 1.215.
Boogbasalt: $3/0.658 = 4.56$, $6.1/0.648 = 9.41$, verhouding 0.484. Alleen de boogbasalt heeft een
niobiumanomalie, een dal bij Nb in een spindiagram: niobium is verarmd ten opzichte van
zijn buur La: het boogmagma kreeg elementen van de plaat, maar weinig
niobium. MORB is bijna vlak; de \omterm{def:b2:elements-in-the-earth:morb-oib}{oceaaneilandbasalt} is rijker
aan Nb dan aan La.
\end{solution}

\begin{solution}{exo:b2:magmatic-series:7}
(a) $\mathrm{Ti} = 1.6\times0.5993\times10^4 = \qty{9589}{ppm}$; Ti/100 = 95.9, Zr
110, $3\mathrm Y = 96$; som 301.9; coördinaten 31.8, 36.4, 31.8: veld B. (b) Ti
16\,780, dus 167.8, 220, 87; som 474.8; 35.3, 46.3, 18.3: veld D. (c) Ti 4794, dus
47.9, 50, 60; som 157.9; 30.3, 31.7, 38.0: veld B. Basalt (b) is een intraplaatbasalt
(van een oceaaneiland of een continent); (a) en (c), in veld B, kunnen MORB,
boogtholeïeten of kalkalkalische basalten zijn: het diagram kan niet beslissen, en een ander (een
spindiagram, de Nb-anomalie) is nodig.
\end{solution}

\begin{solution}{exo:b2:magmatic-series:8}
(a) MORB $142/164 = 0.866$, \omterm{def:b2:elements-in-the-earth:morb-oib}{oceaaneilandbasalten} $4/173 = 0.023$, boogbasalten $101/201 =
0.502$. (b) Met gelijke a-priorikansen geeft de regel van Bayes $0.866/(0.866 + 0.023 + 0.502) =
0.866/1.391 = 0.623$. (c) Dat de basalt zeer waarschijnlijk geen \omterm{def:b2:elements-in-the-earth:morb-oib}{oceaaneilandbasalt} is
(kans $0.023/1.391 = 0.017$), maar niet of hij een MORB of een boogbasalt is:
ruwweg drie kansen op vijf tegen twee.
\end{solution}

\begin{solution}{exo:b2:magmatic-series:9}
Mol per \qty{100}{g}: \ce{Na2O} $5.5/61.98 = 0.08874$, \ce{K2O} $4.5/94.20 =
0.04777$, \ce{Al2O3} $10.5/101.96 = 0.10298$. (a) $(0.08874 + 0.04777)/0.10298 =
0.13651/0.10298 = 1.326 > 1$: de rhyoliet is \emph{peralkalisch}. (b) Veldspaat kan
$0.10298$ van de $0.13651$ mol alkalioxide opnemen; het overschot, $0.03353$, is
$0.03353/0.13651 = 0.246$, een kwart van de alkaliën, dat in andere mineralen moet gaan.
\end{solution}

\begin{solution}{exo:b2:magmatic-series:10}
(a) $s = (\lambda_F - \lambda_M)/(\lambda_F - \lambda_A)$: voor $\lambda_F = 0.5$, $(0.5
- 0.8)/(0.5 - 0.4) = -3.0$; voor $0.7$, $(-0.1)/(0.3) = -0.333$. (b) Oplossen van
$s(\lambda_F - 0.4) = \lambda_F - 0.8$ geeft $\lambda_F = (0.4s - 0.8)/(s - 1)$:
IJsland, $s = -1.10$: $(-0.44 - 0.8)/(-2.10) = 0.590$; Cascaden, $s = -0.72$:
$(-0.288 - 0.8)/(-1.72) = 0.633$. (c) Het \omterm{def:b2:magma-emplacement:magma}{magma} van de Cascaden, met de grotere
$\lambda_F$, verliest zijn ijzer sneller: magnetiet voert het af. (d) $\ln(F/M)$ groeit met de
snelheid $\lambda_M - \lambda_F$ en $\ln(A/F)$ met de snelheid $\lambda_F - \lambda_A$; dus
voor elke eenheid die $\ln(A/F)$ wint, wint $\ln(F/M)$ $-s$: 1.10 in IJsland, 0.72 in de
Cascaden. $F/M$ stijgt langs beide trends, maar sneller langs de tholeïtische, de
Fennertrend.
\end{solution}

\begin{solution}{exo:b2:magmatic-series:11}
(a) BHVO-2, mol per \qty{100}{g}: \ce{SiO2} 0.82557, \ce{Al2O3} 0.13182, \ce{FeO}
0.15521, \ce{MgO} 0.18007, \ce{CaO} 0.20328, \ce{Na2O} 0.03580, \ce{K2O} 0.00545.
Orthoklaas 0.00545 en albiet 0.03580; overgebleven \ce{Al2O3} $0.13182 - 0.04125 = 0.09057$,
geheel in anorthiet met evenveel \ce{CaO}; overgebleven \ce{CaO} $0.11271$ in diopsied;
overgebleven $(\ce{Mg,Fe})\ce{O}$ $0.33528 - 0.11271 = 0.22257$ in hypersteen. Benodigd
silica: $6\times0.00545 + 6\times0.03580 + 2\times0.09057 + 2\times0.11271 + 0.22257 =
0.87663$; beschikbaar 0.82557; tekort 0.05106. Elk $(\ce{Mg,Fe})\ce{O}$ dat van
hypersteen in olivijn wordt omgezet, spaart een half silica, hoogstens $0.22257/2 = 0.11129$: omzetten van
$0.10212$ ervan volstaat, en $0.12045$ blijft in hypersteen. De norm bevat olivijn
en hypersteen, en kwarts noch nefelien: het gesteente is verzadigd aan silica in de
zin van de definitie, een \emph{olivijntholeïet} met hypersteen in de norm, van de
\omterm{def:b2:magmatic-series:series}{tholeïtische reeks}, zoals de schildfase van een Hawaïaanse vulkaan die voortbrengt. (b) Mol: 0.75732, 0.14712, 0.16008, 0.21092, 0.17832, 0.06131,
0.01592. Or 0.01592, ab 0.06131, an 0.06988, di 0.10843, hy $0.37100 - 0.10843 =
0.26257$; benodigd silica $0.09552 + 0.36786 + 0.13976 + 0.21686 + 0.26257 = 1.08257$;
beschikbaar 0.75732; tekort 0.32525. Alle hypersteen in olivijn omzetten spaart
0.13129, wat 0.19396 overlaat; elke mol \ce{Na2O} die van albiet in nefelien wordt omgezet,
maakt 4 \ce{SiO2} vrij, dus $0.19396/4 = 0.04849$
van de 0.06131 mol \ce{Na2O} van het albiet wordt nefelien. Nefelien in de norm: het
gesteente is onderverzadigd aan silica, een \emph{alkalibasalt}, van de \omterm{def:b2:magmatic-series:series}{alkalische reeks}.
\end{solution}

\begin{solution}{exo:b2:magmatic-series:12}
(a) Boogbasalten $116/141 = 0.823$, MORB $21/101 = 0.208$, \omterm{def:b2:elements-in-the-earth:morb-oib}{oceaaneilandbasalten}
$3/137 = 0.022$. (b) $0.823/(0.823 + 0.208 + 0.022) = 0.823/1.053 = 0.782$.
(c) Teller $0.1\times0.823 = 0.0823$; noemer $0.0823 + 0.6\times0.208 +
0.3\times0.022 = 0.0823 + 0.1248 + 0.0066 = 0.2137$; kans $0.385$. (d) Dezelfde
test op hetzelfde gesteente geeft 0.78 of 0.39, afhankelijk van hoe vaak elke omgeving
voorkomt waar het gesteente werd gevonden: een \omterm{def:b2:magmatic-series:discrimination}{discriminatiediagram} geeft een waarschijnlijkheid, en de
geologische context moet de rest leveren.
\end{solution}

\begin{solution}{pb:b2:magmatic-series:1}
\noindent\textbf{Deel I.}

\textbf{1.} Oxiden: lava 1 99.17, lava 2 98.65, lava 3 98.20; met de vluchtige bestanddelen
99.97, 99.80 en 99.45, alle dicht bij 100. Water (0.70, 1.10, 1.20) ligt onder
\qty{2}{\percent} en koolstofdioxide (0.10, 0.05, 0.05) onder \qty{0.5}{\percent}: alle
drie zijn aanvaardbaar.

\textbf{2.} Factor $100/99.17 = 1.00837$: \ce{SiO2} 50.82, \ce{Na2O} 2.92, \ce{K2O}
0.63, alkaliën 3.55.

\textbf{3.} Lava 2: $100/98.65 = 1.01368$, \ce{SiO2} 59.40, \ce{Na2O} 3.95, \ce{K2O}
1.17, alkaliën 5.12. Lava 3: $100/98.20 = 1.01833$, \ce{SiO2} 66.50, \ce{Na2O} 4.48,
\ce{K2O} 1.88, alkaliën 6.36.

\textbf{4.} Lava 1 (50.82, 3.55): basalt. Lava 2 (59.40, 5.12; de bovengrens van O2
ligt daar op $5.9 + 1.1\times2.40/6 = 6.34$): andesiet. Lava 3 (66.50, 6.36; tussen 63
en de lijn van (69, 8) naar (77, 0), onder de lijn van (63, 7) naar (69, 8)): daciet.

\textbf{5.} Onherrekend: (50.4, 3.52), (58.6, 5.05), (65.3, 6.25): dezelfde velden. De
herberekening verschuift elk punt hoogstens \qty{1.2}{\percent} silica en verandert
hier niets, anders dan in \cref{exo:b2:magmatic-series:2}.

\textbf{6.} Het \omterm{def:b2:magmatic-series:qapf}{QAPF-diagram}, uit de gemeten modus. Een plutonisch gesteente met andesitische
samenstelling is rijk aan plagioklaas, met weinig alkaliveldspaat of kwarts: het zou
in veld 10 of 10* moeten vallen, een dioriet of kwartsdioriet (zijn modus beslist).

\noindent\textbf{Deel II.}

\textbf{7.} Lava 1: lijnen op 0.43 en 1.47, $\ce{K2O} = 0.63$: middel-K. Lava 2: 0.81
en 2.28, 1.17: middel-K. Lava 3: 1.13 en 2.96, 1.88: middel-K.

\textbf{8.} \ce{FeO_T} en \ce{MgO} zonder vluchtige bestanddelen: 8.67 en 7.97; 6.18 en 3.85; 4.07
en 1.63. Lava 1: $3.55 + 8.67 + 7.97 = 20.19$; $A = 17.6$, $F = 42.9$, $M = 39.5$.
Lava 2: 15.15; 33.8, 40.8, 25.4. Lava 3: 12.06; 52.7, 33.7, 13.5.

\textbf{9.} Lava 1, $A = 17.6$: de lijn tussen (42.8, 17.1) en (46.9, 18.1) ligt op
$F = 42.8 + 4.1\times0.5/1.0 = 44.9$; $42.9 < 44.9$, kalkalkalische kant, dicht bij de
lijn. Lava 2, $A = 33.8$: $F = 49.6 + 2.9\times1.2/5 = 50.3$; $40.8 < 50.3$,
kalkalkalisch. Lava 3, $A = 52.7$: $F = 40.8 - 2.9\times2.7/5 = 39.2$; $33.7 < 39.2$,
kalkalkalisch.

\textbf{10.} $\ce{FeO_T}/\ce{MgO}$: $8.67/7.97 = 1.09$, $6.18/3.85 = 1.61$, $4.07/1.63 =
2.50$. De verhouding stijgt, zoals in elk differentiërend basalt, maar het ijzer zelf
daalt van 8.67 tot 4.07, en $F$ daalt in de driehoek van 42.9 tot 33.7: geen
ijzerverrijking. Een \omterm{def:b2:magmatic-series:series}{tholeïtische reeks} zou eerst naar F klimmen, zoals de ferrobasalt van
\cref{exo:b2:magmatic-series:5} doet.

\textbf{11.} De \omterm{def:b2:magmatic-series:series}{kalkalkalische reeks}. Magnetiet, gekristalliseerd uit een geoxideerd \omterm{def:b2:magma-emplacement:magma}{magma},
voerde ijzer samen met magnesium af; de geoxideerde species komen van de ontwaterende plaat van een
subductiezone.

\noindent\textbf{Deel III.}

\textbf{12.} $\mathrm{Nb_N} = 2.0/0.658 = 3.04$, $\mathrm{La_N} = 7/0.648 = 10.80$,
$\mathrm{Th_N} = 1.1/0.0795 = 13.84$.

\textbf{13.} $3.04/10.80 = 0.281$, ruim onder de medianen van MORB (0.956) en
\omterm{def:b2:elements-in-the-earth:morb-oib}{oceaaneilandbasalten} (1.215), en zelfs onder de mediane boogbasalt (0.484): een diep
niobiumdal.

\textbf{14.} $13.84/3.04 = 4.55$; de mediane MORB van het spindiagram heeft $2.64/6.08 =
0.43$. Thorium werd, net als de andere elementen die in boogbasalten verrijkt zijn, door de
subducerende plaat toegevoegd.

\textbf{15.} \ce{TiO2} zonder vluchtige bestanddelen $0.95\times1.00837 = 0.958$; $\mathrm{Ti} =
0.958\times5993 = \qty{5741}{ppm}$; Ti/100 = 57.4, Zr 70, $3\mathrm Y = 54$, som 181.4:
31.6, 38.6, 29.8, in veld B, gedeeld door MORB, boogtholeïeten en kalkalkalische basalten.
Het sluit een intraplaatbasalt uit, maar beslist niet tussen een rug en een boog. Het
diagram is gemaakt voor basalten: in een andesiet en een daciet is titanium verwijderd
door de oxiden en zirkonium geconcentreerd door kristallisatie, en hun velden (lava 2
zou in C vallen, lava 3 buiten elk veld) zouden niets betekenen.

\textbf{16.} Een kalkalkalische middel-K-reeks, met een diep niobiumdal en een hoge
$\mathrm{Th_N}/\mathrm{Nb_N}$, op een vulkaan die basalt tot daciet voortbracht: de vulkaan
groeide in een vulkanische boog boven een subductiezone.

\noindent\textbf{Deel IV.}

\textbf{17.} $C_L = C_0/F$ voor een volmaakt \omterm{def:b2:elements-in-the-earth:partition}{incompatibel element}: $F = 0.63/1.17 =
0.538$.

\textbf{18.} Rubidium $10/20 = 0.500$; thorium $1.1/2.1 = 0.524$. (De
\omterm{def:b2:elements-in-the-earth:trace}{sporenelementen}, gegeven zoals geanalyseerd, zouden door de herberekening minder dan \qty{1}{\percent}
veranderen.)

\textbf{19.} $C_L/C_0 = F^{D - 1}$: $125/70 = 1.786 = F^{-0.9}$, $F = 1.786^{-1/0.9} =
0.525$.

\textbf{20.} Uit kali $F = 1.17/1.88 = 0.622$ van lava 2 naar lava 3. Strontium:
$380/480 = 0.792 = 0.622^{D - 1}$, dus $D - 1 = \ln 0.792/\ln 0.622 = -0.234/-0.474 =
0.49$ en $D \approx 1.49$: strontium is compatibel geworden. Plagioklaas, dat
strontium opneemt, kristalliseerde in overvloed.

\textbf{21.} Uit kali $F = 0.63/1.88 = 0.335$ (rubidium $10/32 = 0.31$, thorium
$1.1/3.4 = 0.32$): ongeveer twee derde van de basalt, $1 - 0.335 = 0.665$, kristalliseerde
vóór de daciet.

\textbf{22.} Geassimileerde kali zou het \ce{K2O} in lava 2 verhogen boven wat
kristallisatie alleen geeft, dus $F = 0.63/1.17$ zou te laag uitvallen, en de massa
kristallen te hoog.

\textbf{23.} $(0.538 + 0.500 + 0.524 + 0.525)/4 = 2.087/4 = 0.522$; kristallen $1 - 0.522
= 0.478$ van de basalt, dat wil zeggen \textbf{\qty[text-series-to-math=true]{48}{kg}
kristallen per \qty{100}{kg} basalt} om de andesiet te maken.
\end{solution}
